\documentclass[11pt,a4paper]{amsart}

\usepackage{amsmath,amssymb,amsthm}
\usepackage[mathscr]{euscript}
\usepackage[T1]{fontenc}
\usepackage{lmodern}
\usepackage[hidelinks]{hyperref}
\usepackage{cleveref}
\usepackage{enumitem}
\usepackage{booktabs,array}
\usepackage[protrusion=true,expansion=false]{microtype}
\usepackage{mathtools}
\raggedbottom
\emergencystretch=3em
\allowdisplaybreaks

\theoremstyle{plain}
\newtheorem{theorem}{Theorem}[section]
\newtheorem{proposition}[theorem]{Proposition}
\newtheorem{corollary}[theorem]{Corollary}
\newtheorem{lemma}[theorem]{Lemma}

\theoremstyle{definition}
\newtheorem{definition}[theorem]{Definition}

\theoremstyle{remark}
\newtheorem{remark}[theorem]{Remark}

\DeclareMathOperator{\Tr}{Tr}
\DeclareMathOperator{\spec}{spec}
\DeclareMathOperator{\Irr}{Irr}
\DeclareMathOperator{\Hom}{Hom}
\DeclareMathOperator{\Aut}{Aut}
\DeclareMathOperator{\End}{End}
\DeclareMathOperator{\Log}{Log}
\DeclareMathOperator{\rad}{rad}
\newcommand{\CC}{\mathbb{C}}
\newcommand{\RR}{\mathbb{R}}
\newcommand{\ZZ}{\mathbb{Z}}
\newcommand{\NN}{\mathbb{N}}
\newcommand{\QQ}{\mathbb{Q}}
\newcommand{\TT}{\mathbb{T}}
\newcommand{\cP}{\mathcal P}
\newcommand{\abs}[1]{\lvert #1 \rvert}
\newcommand{\norm}[1]{\lVert #1 \rVert}
\newcolumntype{L}[1]{>{\raggedright\arraybackslash}p{#1}}

\title[Conditional finite radial determination]{Finite radial determination
under a Mahler rationality hypothesis for odd-Adams-invariant abelian
two-group extensions of shifts of finite type}

\author[Haobo Ma]{Haobo Ma}
\address{AELF PTE. LTD., \#14-02, Marina Bay Financial Centre Tower 1,
8 Marina Boulevard, Singapore 018981, Singapore}
\email{auric@aelf.io}
\author[Wenlin Zhang]{Wenlin Zhang}
\address{National University of Singapore, 21 Lower Kent Ridge Road,
Singapore 119077, Singapore}
\email{e1327962@u.nus.edu}
\thanks{Corresponding author: Wenlin Zhang.}
\date{}
\renewcommand{\shortauthors}{H. Ma and W. Zhang}

\begin{document}

\begin{abstract}
Let \(G\) be a finite abelian \(2\)-group, and assume that the twisted
determinants of two one-step \(G\)-cocycles are invariant under every odd
Adams operation \(\chi\mapsto\chi^u\).  The cocycles lie over primitive edge
shifts whose vertex counts are at most \(V\); the bases need not be the same.
For the primitive length--element counts \(p_{n,g}\), put
\(\mathcal E_g(y)=\sum_{n\ge1}p_{n,g}\log(1-y^n)\).  Assume the named
statement \textup{(KN85)} quoted in the paper.  Without a twisted-gap
hypothesis, equality of every element profile at
\[
 2V\lceil\log_2(4V)\rceil
\]
distinct radii in the common open Perron interval forces equality of every
\(p_{n,g}\), provided that one sampled radius is algebraic.  Thus the number
of radial locations is independent of the rank and exponent of \(G\), although
each location supplies the full \(\abs G\)-component profile vector.  The
hypothesis is automatic for \(G=(C_2)^r\) and permits genuine holonomy of
orders \(4,8,\ldots\).  For a reduced binary character-determinant ratio of
total degree \(D>0\), one algebraic collision identifies the complete real
collision set, which has at most
\(D\lceil\log_2(2D)\rceil-1\) points.

The sign-determinant ratio is canonically the ratio of the ordinary
Artin--Mazur zeta functions of the two \(C_2\)-covers.  The inverse mechanism
is arithmetic: a single algebraic interior collision forces this ratio to be
a normalized rational Mahler coboundary
\(R(z^2)/R(z)^2\).  More generally, a rational critical \(p\)-Mahler product
with one algebraic value is algebraic by an elementary linear-exponent
denominator estimate and Kumiko Nishioka's critical special-value theorem.
The finite-sampling and algebraic-collision upper bounds rest on
\textup{(KN85)}, a named statement attributed to a published 1985 paper whose
bibliographic record is verified but whose printed theorem the authors were
unable to consult; the exact statement required is quoted in the paper.  All
results not explicitly carrying \textup{(KN85)} are unconditional.  In
particular, the sharp squarefree Mahler bound, the collision--jet inequality,
the odd-prime multi-collision theorem, and the sampling lower bound require no
such input.  Determinant parity and the
Dieudonn\'e--Dwork criterion give a stronger integral refinement but are not
needed for lifting.  For the multiplicative certificate \(R=A/B\), a new
squarefree divisor estimate gives
\[
 2(p-1)\deg\rad(AB)\le D.
\]
Together with Rolle's theorem, it bounds the number of positive collisions
plus the first non-zero Taylor order of \(R-1\).  A separate divisor
\(\ell_1\) estimate gives the input-only total-degree bound
\(\deg A+\deg B\le b_p(D)=O_p(D\log D)\) in terms of the reduced input degree
\(D\).  An \(\Omega_p(D\log D)\) family makes this order sharp, and the same
order occurs for standard realizable \(C_2\)-cover zeta ratios.  Explicit
height control and fixed-\(p\) polynomial bit complexity are further
quantitative consequences; once the degree cap is known, the displayed
Pad\'e step is rational reconstruction.  Conditional on \textup{(KN85)}, the
squarefree estimate gives a linear prime-primary collision--jet theorem.
An exact four-vertex pair shows that one centered
Perron-boundary profile vector does not suffice; the uncentered class profiles
diverge there.  More generally, the binary construction on a common
\(4m+2\)-vertex positive base gives \(m\) rational collisions, while for
every odd prime \(\ell\) a pair on \(\ell(2m+1)\) vertices has \(m\) rational
collisions, agrees through length \(m\), and first differs explicitly at
length \(m+1\).  A companion-matrix realization transfers the
\(\Omega(D\log D)\) certificate-degree family to standard zeta ratios.  For general finite
groups we prove
recovery from a radial set with an interior accumulation point and identify
the obstruction to a finite bound: already for \(C_3\), Adams--M\"obius
inversion couples determinant logarithms at infinitely many arithmetically
distributed powers, outside the one-variable Mahler theorem used here.
\end{abstract}

\subjclass[2020]{37B10, 11B85, 39A06, 20C15, 05C25}
\keywords{One-sided edge shift, finite-group extension, inverse problem,
Mahler function, twisted determinant, Adams operation, primitive orbit,
finite sampling}

\maketitle
\enlargethispage{4pt}

\section{Introduction}
\label{sec:introduction}

Let \(X_\Gamma^+\) be the one-sided edge shift of a finite primitive directed
multigraph, and let a one-step cocycle \(\tau:E(\Gamma)\to G\) take values in a
finite group.  A primitive orbit \(\gamma\) has a well-defined conjugacy class
\([g_\gamma]\).  Write \(p_{n,C}(\tau)\) for the number of primitive
length-\(n\) orbits with class \(C\), and define its radial class Euler profile
by
\begin{equation}\label{eq:intro-class-euler-profile}
 \mathcal E_C^\tau(y)
 :=\sum_{n\ge1}p_{n,C}(\tau)\log(1-y^n).
\end{equation}
This converges without a twisted-gap hypothesis for
\(0<y<\lambda^{-1}\), where \(\lambda\) is the Perron root of the base.
Under the strict twisted gap, every raw class profile diverges to \(-\infty\)
at \(y=\lambda^{-1}\), and no finite endpoint value of \(\mathcal E_C\) is
used.  The finite endpoint coordinate used below is the
centered profile
\[
 \mathcal E_C^{\circ}(y)
 :=\mathcal E_C(y)-\frac{|C|}{|G|}\log\det(I-yA),
 \qquad 0<y<\lambda^{-1}.
\]
Under the strict gap its radial limit exists by
\Cref{prop:centered-perron-profile}.  The inverse question of this paper is
whether finitely many open-interval values of all the raw profiles determine
every \(p_{n,C}\); endpoint statements concern only the centered profiles or
their equivalent non-trivial Fourier coordinates.

The determinant input is finite.  For each irreducible representation \(\rho\)
of \(G\), the twisted adjacency block \(A_\rho\) has size
\(v\dim\rho\), where \(v=|V(\Gamma)|\), and
\(P_\rho(z)=\det(I-zA_\rho)\).  What makes the radial inverse problem
non-trivial is that \(\mathcal E_C\) is not the logarithm of one of these
determinants.  For a class function \(\phi\), the periodic Artin logarithm is
\begin{equation}\label{eq:intro-artin}
 L_\phi(z)=\sum_{\gamma\in\mathcal P}\sum_{r\ge1}
   \frac{\phi(g_\gamma^r)}r z^{r|\gamma|},
\end{equation}
whereas the fixed-label coordinate is
\begin{equation}\label{eq:intro-fixed}
 F_\phi(z)=\sum_{\gamma\in\mathcal P}\sum_{r\ge1}
   \frac{\phi(g_\gamma)}r z^{r|\gamma|}.
\end{equation}
Thus \(\mathcal E_C^\tau=-F_{\mathbf1_C}^\tau\).  Adams--M\"obius
inversion expresses every \(F_\phi\) through the finite determinant family,
but it does so at all powers \(z^k\).  Finite radial determination is therefore
a special-value problem, not a coefficient-counting consequence of the
determinant degrees.

\subsection{Main results}

We first state the sharper scalar estimate for binary extensions on a fixed base.  Let
\(G=C_2\), let \(B,B'\) be the two sign blocks, and reduce
\[
 H(z)=\frac{\det(I-zB)}{\det(I-zB')}=\frac{P_0(z)}{P_1(z)}
\]
with \(D=\deg P_0+\deg P_1\).  By
\Cref{prop:c2-standard-zeta-ratio}, \(H\) is the ratio of the ordinary
Artin--Mazur zeta functions of the two standard binary covers.  Without any
twisted-gap assumption, if \(H\ne1\), then
\Cref{thm:finite-radial-sampling} proves
\begin{equation}\label{eq:intro-collision-bound}
 \#\{y\in\overline{\mathbb Q}\cap(0,\lambda^{-1}):
   \mathcal E_C^\tau(y)=\mathcal E_C^{\tau'}(y)
   \text{ for every }C\}
 \le D\lceil\log_2(2D)\rceil-1.
\end{equation}
Consequently equality of every class profile at any
\begin{equation}\label{eq:intro-universal-sample-budget}
 B(v,C_2)=2v\lceil\log_2(4v)\rceil
\end{equation}
distinct radii in the open Perron interval, at least one of which is algebraic,
forces
\(p_{n,C}(\tau)=p_{n,C}(\tau')\) for all \(n,C\).  The radii need not be
chosen adaptively.  If both sign blocks have spectral radius below
\(\lambda\), the collision and recovery statements extend to the Perron
endpoint for centered class-profile differences, not for the divergent raw
profiles.

The proof has two independent parts.  First,
\Cref{thm:determinant-boundary-lifting} shows that one algebraic interior
collision, without a spectral gap, or one endpoint collision under the gap,
forces the exact rational Mahler coboundary
\[
 H(z)=\frac{R(z^2)}{R(z)^2},\qquad R(0)=R(y)=1.
\]
This is the delicate step.  Determinant parity makes the associated critical
Mahler product integral; Kumiko Nishioka's special-value theorem then forces it
to be algebraic, and Keiji Nishioka's 1985 algebraic-solution theorem makes it rational.
We verify every
hypothesis of Kumiko Nishioka's printed theorem in the proof, including the critical
numerical inequality.  Second,
\Cref{thm:effective-rational-mahler-coboundary} bounds the reduced numerator
and denominator degrees of \(R\) and reconstructs it by one finite Pad\'e
system.  The logarithmic divisor estimate in
\Cref{prop:logarithmic-mahler-divisor-bound} has optimal order at the level of
general rational divisor equations.  No matching lower bound is claimed for
realizable SFT collision sets.  The collision radii are exactly the zeros of
\(R-1\), giving
\eqref{eq:intro-collision-bound}.

The principal inverse theorem is not confined to a common base or to one
binary coordinate.  By \Cref{thm:elementary-two-group-finite-sampling}, for
\(E=(C_2)^r\), two primitive bases of at most \(V\) vertices are determined at
the level of all primitive length--element counts by
\[
 M(V)=2V\lceil\log_2(4V)\rceil
\]
common profile samples, with one algebraic anchor.  Summing the element
profiles first recovers the base determinant.  Reduction modulo two then
supplies the critical integrality input for each of the \(2^r\) character
blocks, while the same scalar zero bound applies separately to every block.
Here one observation at a radius means the entire \(E\)-valued vector
\((\mathcal E_g(y))_{g\in E}\).  Hence the theorem uses \(M(V)\)
radial locations independently of \(r\), but it compares
\(2^rM(V)\) scalar profile values.

One centered radius cannot suffice.  \Cref{thm:c2-boundary-collision} gives
two strict-gap labellings of the same four-vertex graph with different
length-one class counts but identical centered class Euler profiles at the
Perron radius.  Their raw class profiles both diverge there.  The centered
identity is exact and comes from
\(Q(z^2)/Q(z)^2\), where \(Q(z)=1-z+4z^2\) and \(Q(1/4)=1\).

For an arbitrary finite group, \Cref{thm:radial-profile-recovery} proves,
again without a twisted-gap hypothesis, recovery from a set with an interior
accumulation point.  We do not claim a
finite polynomial sample bound in that generality.  The precise obstruction
to extending the present method is recorded after
\Cref{thm:finite-radial-sampling}: already for \(C_3\), the
Adams--M\"obius coefficient is non-zero at infinitely many prime powers and
does not reduce to a one-variable Mahler orbit.  The missing input is a
special-value and effective zero theorem for that coupled cyclotomic Euler
system.  This locates the failure of the proposed general-group tier-up
statement rather than replacing it by a weaker finite-length theorem.

\subsection{Prior work and scope}

The Frobenius-class objects used here have a direct dynamical lineage.
Parry--Pollicott prove a Chebotarev asymptotic for primitive closed orbits in a
prescribed Frobenius conjugacy class in finite Galois covers of Axiom~A flows,
and Noorani--Parry prove the finite homogeneous-extension analogue for shifts
\cite{ParryPollicott1986Chebotarev,NooraniParry1992ChebotarevShifts}.
Mohamed--Noorani then prove a class-restricted Mertens product for finite-group
extensions of SFTs: their Theorem~1 has the product
\[
 \prod_{\substack{|\gamma|\le N\\g_\gamma\in C}}
 (1-e^{-h|\gamma|})
\]
with exponent \(|C|/|G|\) and a constant expressed through Artin
\(L\)-values \cite[Theorem~1, pp.~124--125]
{MeslizaNoorani1999FrobeniusMertens}.  Noorani's later toral paper instead
classifies the lift-degree partitions of closed orbits in a finite torsion
cover and gives asymptotic counts for each realizable partition
\cite{Noorani2003ClosedOrbitsToral}.  These are orbit-distribution and product
results, not finite-special-value inverse theorems.

At the formula level, the standard Artin function in that literature has
\[
 \log L(z,\phi)
 =\sum_{\gamma}\sum_{m\ge1}
   \frac{\phi(g_\gamma^m)}m z^{m|\gamma|}
 =L_\phi(z),
\]
whereas taking the logarithm of a product whose membership condition remains
\([g_\gamma]=C\) gives
\[
 \sum_{[g_\gamma]=C}\log(1-z^{|\gamma|})
 =-F_{\mathbf1_C}(z),\qquad
 F_\phi(z)=\sum_{\gamma}\sum_{m\ge1}
   \frac{\phi(g_\gamma)}m z^{m|\gamma|}.
\]
Thus class-restricted Euler products, character orthogonality, and the
\(\Pi_\phi,L_\phi\) coordinates are prior art.  The fixed-label transform
\(F_\phi\) is already implicit in the logarithm of the 1999 product.  What is
specific here is that \(F_\phi\), rather than \(L_\phi\), is passed through the
Adams--M\"obius formula \eqref{eq:fixed-adams-mobius} and then used in a
finite-special-value inverse problem.
\begin{center}
\small
\begin{tabular}{@{}L{0.27\textwidth}L{0.29\textwidth}L{0.36\textwidth}@{}}
\toprule
Coordinate or assertion & Prior status & Status in this paper \\
\midrule
Frobenius-class orbit count & Chebotarev asymptotics in
\cite{ParryPollicott1986Chebotarev,NooraniParry1992ChebotarevShifts}
& Used as background; no new equidistribution theorem. \\
\(\Pi_\phi\) and character projection & Class restriction and character
orthogonality occur in the cited Chebotarev--Artin framework
& Standard primitive coordinate; no originality claim. \\
\(L_\phi\) & Logarithm of the standard Artin \(L\)-function
& Standard determinant coordinate. \\
\(-F_{\mathbf1_C}\) & Logarithm of the fixed-class Euler product in
\cite{MeslizaNoorani1999FrobeniusMertens}
& Adams--M\"obius determinant readout and finite radial inversion are the
paper-specific steps. \\
\bottomrule
\end{tabular}
\end{center}

The algebraic-solution rationality step is also prior.  Ostrowski treats the
linear multiplicative equation
\(\Phi(\varphi(z))=g(z)\Phi(z)\) for rational data
\cite{Ostrowski1968AlgebraicSolutions}; that equation does not directly cover
the square in \(F(z^2)=H(z)^{-1}F(z)^2\).  Keiji Nishioka's 1985 theorem does:
its one-variable specialization says that if
\(f(z^p)=\mathscr R(z,f(z))\), with \(\mathscr R\) rational and \(f\)
algebraic over \(\CC(z)\), then \(f\in\CC(z)\)
\cite{Nishioka1985AlgebraicSolutions}.  Taking
\(p=2\) and \(\mathscr R(z,Y)=Y^2/H(z)\) proves the rationality used in
\Cref{prop:algebraic-mahler-coboundary}; that rationality is not claimed as a
contribution of this paper.

The nearest algorithms compute rational solutions of linear Mahler equations
\cite{ChyzakDreyfusDumasMezzarobba2018MahlerSolutions} and first-order
factors through Riccati equations whose monomials are products of successive
shifts of the unknown
\cite{ChyzakDreyfusDumasMezzarobba2025FirstOrderFactors}.  The present equation
\(P_0(z)R(z)^2=P_1(z)R(z^2)\) repeats the unshifted unknown; a rational
monomial substitution preserves that repeated exponent and does not turn it
into a product of distinct successive shifts.  We therefore do not claim the
first decision procedure for general Mahler--Riccati problems or a lower
complexity than those algorithms.  The effective addition here is the
input-only \(O(D\log D)\) degree bound, an explicit height bound, and a single
finite Pad\'e reconstruction for this normalized multiplicative equation.
Additive Mahler summability \cite{ArrecheZhang2022MahlerResidues} concerns a
different equation.

Accordingly, the claimed contributions are limited to three linked steps:
parity-compatible algebraic collision lifting (with Kumiko Nishioka's
special-value theorem and Keiji Nishioka's algebraic-solution rationality
theorem cited for their respective steps), effective rational-coboundary bounds and
reconstruction, and cross-base finite radial determination for elementary
two-groups.  The class Euler products, their character expansion, the general
periodic-data dictionary, Keiji Nishioka's algebraic-to-rational theorem, and the
analytic identity-theorem recovery statement are background rather than
novelty claims.  Boyle--Schmieding's computable periodic invariants
\cite{BoyleSchmieding2017FiniteGroupExtensions} and represented
cospectrality and switching in gain-graph theory
\cite{Zaslavsky1989BiasedGraphs,CavaleriDAngeliDonno2021BalanceGainGraphs,
CavaleriDonno2022Cospectrality} delimit the dynamical side further.

Three printed open-problem interfaces delimit the claim.  Boyle and
Schmieding ask to ``characterize the polynomials \(\det(I-tA)\) arising from
\(G\)-primitive matrices'' and to ``characterize the trace series \(T_A\) and
conjugate trace series \(\kappa T_A\)''
\cite[Realization Problems~6.1(2)--(3)]{BoyleSchmieding2017FiniteGroupExtensions}.
Our theorem starts with two realized extensions, so its determinant and trace
formulas do not solve either positivity-realization problem.  Chyzak,
Dreyfus, Dumas, and Mezzarobba leave open polynomial-time solution of general
linear Mahler equations for all order--degree regimes
\cite[Sec.~4.3]{ChyzakDreyfusDumasMezzarobba2018MahlerSolutions}; our finite
Pad\'e criterion treats one normalized nonlinear coboundary equation and does
not settle that complexity problem.  O'Hare asks for finite-data rigidity for
unimodal and Markov maps \cite[Introduction]{OHare2025FiniteDataRigidity}.
The present theorem is an exact symbolic result for orbit-class Euler data,
not a derivative-data rigidity theorem for those maps.

The separately submitted Supplementary Material contains the Frobenius-class Mertens correction from which the profiles arose, its
quotient-cover refinement, and the exact \(S_3\) witness are now separated
from the inverse argument.  This paper uses only
the fixed-label coordinate and does not claim that the underlying Mertens
product or Frobenius equidistribution is new.

\subsection{Organization}

Section~\ref{sec:framework} establishes the represented determinant and
Adams--M\"obius formulas.  Section~\ref{sec:boundary-collisions} contains the
collision, lifting, effective coboundary, finite-sampling, and general radial
recovery theorems.  Section~\ref{sec:exactness} records the verification
boundary.  The appendices give the radial continuation convention and the
standard periodic-data dictionary.

\section{Finite-state framework and the two logarithms}\label{sec:framework}

\subsection{One-sided edge shifts and twisted determinants}

Throughout, \(\Gamma=(V,E,o,t)\) is a finite directed multigraph.  Parallel
edges are allowed and are regarded as named edges.  Its adjacency matrix is
\[
  A_{ij}=\#\{e\in E:o(e)=i,\ t(e)=j\}.
\]
We assume that \(A\) is primitive and that its Perron root satisfies
\(\lambda>1\).  The phase space in every main statement is the one-sided
edge shift
\[
  X_\Gamma^+=\{(e_0,e_1,\ldots):t(e_j)=o(e_{j+1})\},
  \qquad \sigma(e_0e_1\ldots)=e_1e_2\ldots.
\]
No assertion about arbitrary two-sided cocycles is implicit in this
notation.

Fix a finite group \(G\) and an edge labelling \(\tau:E\to G\).  For a
directed path \(q=e_0\cdots e_{n-1}\), set
\(\tau(q)=\tau(e_0)\cdots\tau(e_{n-1})\).  Changing the initial point on a
closed path conjugates this product.  Hence a periodic orbit \(\gamma\) has a
well-defined Frobenius conjugacy class \([g_\gamma]\), although a representative
\(g_\gamma\) requires a base point.

For a finite-dimensional unitary representation
\(\rho:G\to U(V_\rho)\), define the twisted adjacency matrix by blocks
\begin{equation}\label{eq:twisted-adjacency}
  (A_\rho)_{ij}:=\sum_{e:i\to j}\rho(\tau(e))
  \quad\in\End(V_\rho),
  \qquad
  P_\rho(z):=\det(I-zA_\rho).
\end{equation}
For the trivial representation, \(A_{\mathbf1}=A\).  Reversing every
representation convention replaces labels by inverses and classes by inverse
classes; none of the distinction between \(F\) and \(L\) depends on that
choice.

Let \(N_{n,C}\) be the number of period-\(n\) points whose length-\(n\)
holonomy lies in the conjugacy class \(C\).  Let \(p_{n,C}\) be the number of
primitive periodic orbits of least period \(n\) with Frobenius class \(C\).
For a class function \(\phi:G\to\CC\), use the abbreviations
\begin{align}
  T_n(\phi)&:=\sum_C\phi(C)N_{n,C},
  &c_n(\phi)&:=\sum_C\phi(C)p_{n,C},\label{eq:weighted-counts}\\
  \Pi_\phi(z)&:=\sum_{n\geq1}c_n(\phi)z^n
  =\sum_{\gamma\in\cP}\phi(g_\gamma)z^{|\gamma|}.
  \label{eq:primitive-channel}
\end{align}

\begin{lemma}[Twisted trace and periodic logarithm]\label{lem:trace-log}
For every irreducible representation \(\rho\) and every \(n\geq1\),
\begin{equation}\label{eq:twisted-trace}
  \Tr(A_\rho^n)=T_n(\chi_\rho).
\end{equation}
Consequently, on \(|z|<\lambda^{-1}\),
\begin{equation}\label{eq:artin-determinant-log}
  L_{\chi_\rho}(z):=-\Log_0P_\rho(z)
  =\sum_{n\geq1}\frac{T_n(\chi_\rho)}{n}z^n
  =\sum_{r\geq1}\frac{\Pi_{\psi^r\chi_\rho}(z^r)}{r}.
\end{equation}
Here \(\Log_0P_\rho\) is the analytic germ at zero normalized by
\(\Log_0P_\rho(0)=0\).
\end{lemma}

\begin{proof}
Expanding the block product in \(A_\rho^n\) sums \(\rho(\tau(q))\) over
closed length-\(n\) paths, with a trace over \(V_\rho\); this is
\eqref{eq:twisted-trace}.  The finite-matrix identity
\[
  -\Log_0\det(I-zB)=\sum_{n\geq1}\Tr(B^n)z^n/n
\]
gives the second expression in \eqref{eq:artin-determinant-log}.  A primitive
orbit \(\gamma\) of length \(d\) contributes its \(d\) cyclic base points to
period \(n=rd\), and its repeated holonomy is \(g_\gamma^r\).  Division by
\(n=rd\) leaves the factor \(1/r\), proving the last expression.  Absolute
convergence follows from the finite-matrix trace bound on every smaller disc.
\end{proof}

For a general class function, write
\[
  \phi=\sum_{\rho\in\Irr(G)}
     \widehat\phi(\rho)\chi_\rho,
  \qquad
  \widehat\phi(\rho)
    =\frac1{|G|}\sum_{g\in G}\phi(g)\overline{\chi_\rho(g)},
\]
and define \(L_\phi=\sum_\rho\widehat\phi(\rho)L_{\chi_\rho}\).  This agrees
with the orbit expansion in \eqref{eq:artin-determinant-log}.  The Adams
operation \(\psi^m\) acts on class functions by
\((\psi^m\phi)(g)=\phi(g^m)\).  It need not preserve irreducibility or
positivity, but it has a finite character expansion
\begin{equation}\label{eq:adams-coefficients}
  \psi^m\chi_\rho=\sum_{\pi\in\Irr(G)}
      a_{\rho,\pi}^{(m)}\chi_\pi,
  \qquad
  a_{\rho,\pi}^{(m)}
     =\langle\psi^m\chi_\rho,\chi_\pi\rangle_G.
\end{equation}
The coefficients are periodic in \(m\) modulo the exponent of \(G\), and
therefore uniformly bounded.

\subsection{Adams--M\"obius primitive inversion}

\begin{theorem}[Primitive and fixed-label determinant coordinates]
\label{thm:adams-mobius}
For every class function \(\phi\) and \(|z|<\lambda^{-1}\),
\begin{align}
  \Pi_\phi(z)
    &=\sum_{m\geq1}\frac{\mu(m)}{m}
       L_{\psi^m\phi}(z^m),
       \label{eq:primitive-adams-mobius}\\
  F_\phi(z):=\sum_{r\geq1}\frac{\Pi_\phi(z^r)}{r}
    &=\sum_{r,m\geq1}\frac{\mu(m)}{rm}
       L_{\psi^m\phi}(z^{rm}).
       \label{eq:fixed-adams-mobius}
\end{align}
For \(\phi=\chi_\rho\), insertion of \eqref{eq:adams-coefficients} gives a
formula involving only the finite determinant family:
\begin{equation}\label{eq:fixed-determinant-open-disc}
  F_{\chi_\rho}(z)=
  \sum_{r,m\geq1}\frac{\mu(m)}{rm}
    \sum_{\pi\in\Irr(G)}a_{\rho,\pi}^{(m)}
      \bigl(-\Log_0P_\pi(z^{rm})\bigr).
\end{equation}
All three series are absolutely convergent on compact subsets of the Perron
disc after expansion into primitive orbits.
\end{theorem}

\begin{proof}
By \Cref{lem:trace-log},
\[
  L_{\psi^m\phi}(z^m)
  =\sum_{\gamma\in\cP}\sum_{r\geq1}
     \frac{\phi(g_\gamma^{mr})}{r}z^{mr|\gamma|}.
\]
Multiplying by \(\mu(m)/m\), summing, and putting \(s=mr\), the coefficient
of \(\phi(g_\gamma^s)z^{s|\gamma|}\) becomes
\[
  \sum_{m\mid s}\frac{\mu(m)}{m(s/m)}
  =\frac1s\sum_{m\mid s}\mu(m).
\]
It is one for \(s=1\) and zero otherwise.  This proves
\eqref{eq:primitive-adams-mobius}.  Summing that identity over \(z^r/r\)
proves \eqref{eq:fixed-adams-mobius}.  Finally use the character expansion
\eqref{eq:adams-coefficients} and \eqref{eq:artin-determinant-log}.

For convergence, the total number \(p_n\) of primitive length-\(n\) orbits
satisfies \(p_n\leq\Tr(A^n)/n\ll\lambda^n/n\).  Thus the orbit expansions
are dominated on \(|z|\leq R<\lambda^{-1}\) by convergent logarithmic series.
This also justifies every rearrangement above.
\end{proof}

The difference between the periodic and fixed-label transforms is now
formal rather than terminological.

\begin{proposition}[Exact Artin-substitution obstruction]
\label{prop:artin-substitution-obstruction}
On \(|z|<\lambda^{-1}\),
\begin{equation}\label{eq:exact-f-minus-l}
  F_\phi(z)-L_\phi(z)
  =\sum_{r\geq2}\frac{
       \Pi_\phi(z^r)-\Pi_{\psi^r\phi}(z^r)}{r}.
\end{equation}
Writing \(\Pi_\phi(z)=\sum_{n\geq1}c_n(\phi)z^n\), the two germs agree if
and only if
\begin{equation}\label{eq:coefficient-obstruction}
  \sum_{r\mid n}\frac1r
    c_{n/r}(\phi-\psi^r\phi)=0
  \qquad(n\geq1).
\end{equation}
There is no general character-theoretic reason for these identities to hold.
\end{proposition}

\begin{proof}
Subtract the two orbit expansions in
\eqref{eq:intro-artin}--\eqref{eq:intro-fixed}; the \(r=1\) terms cancel.
Extracting the coefficient of \(z^n\) gives
\eqref{eq:coefficient-obstruction}.
\end{proof}

\subsection{The Perron boundary is derived from the determinant}

Put
\begin{equation}\label{eq:strict-gap}
  \eta:=\max_{\rho\in\Irr(G),\ \rho\neq\mathbf1}\rad(A_\rho).
\end{equation}
The main product theorem assumes \(\eta<\lambda\).  Under this hypothesis,
the \(rm=1\) terms in \eqref{eq:fixed-determinant-open-disc} are regular for
non-trivial \(\pi\), and every term with \(rm\geq2\) is an interior value
because
\begin{equation}\label{eq:interior-radius}
  \rad(\lambda^{-rm}A_\pi)\leq\lambda^{1-rm}<1.
\end{equation}

We record the boundary convention in a form that makes clear that no branch
or peripheral factor is an additional input coordinate.

\begin{proposition}[Canonical radial boundary construction]
\label{prop:derived-boundary}
Let \(P_\pi(z)=\det(I-zA_\pi)\).  The normalized germ
\(-\Log_0P_\pi(z)\), its continuation along \(z=t/\lambda\) with
\(0\leq t<1\), its peripheral factorization, and every radial renormalized
endpoint are determined by \(P_\pi\), \(\lambda\), and this path convention.
More explicitly, factor
\begin{equation}\label{eq:peripheral-factorization}
  P_\pi(t/\lambda)
   =\prod_{\zeta\in\Omega_\pi}(1-t\zeta)^{m_{\pi,\zeta}}
      Q_{\partial,\pi}(t),
\end{equation}
where \(\Omega_\pi\subset\TT\) is the multiset of peripheral phases and
\(Q_{\partial,\pi}\) has no zero on \([0,1]\).  If
\(s_\pi=m_{\pi,1}\), then
\begin{align}
  B_{1,\pi}^{\mathrm{ren}}
  &:={\lim}_{t\uparrow1}
     \left(-\Log_{[0,t]}P_\pi(t/\lambda)+s_\pi\log(1-t)\right)
     \label{eq:derived-b-ren}\\
  &=-\sum_{\substack{\zeta\in\Omega_\pi\\\zeta\neq1}}
       m_{\pi,\zeta}\Log_{[0,1]}(1-\zeta)
     -\Log_{[0,1]}Q_{\partial,\pi}(1).
     \label{eq:derived-b-factor}
\end{align}
In particular, a list of branch rules, peripheral factors, or values such as
\(B_{k,\pi}^{\mathrm{ren}}\) is derived output, not independent cocycle data.
\end{proposition}

\begin{proof}
The roots and their multiplicities in \eqref{eq:peripheral-factorization}
are determined by the polynomial \(P_\pi(t/\lambda)\).  The normalized
logarithm has a unique continuation along any zero-free subinterval beginning
at zero.  Removing the factor \((1-t)^{s_\pi}\), taking the radial limit, and
using the same continuation on every remaining factor yields
\eqref{eq:derived-b-factor}.  The construction for \(k>1\) is the identical
one applied at \(t^k/\lambda^k\); by \eqref{eq:interior-radius} it is an
ordinary interior logarithm in the strict-gap theorem below.
\end{proof}

\begin{lemma}[Perron evaluation of a non-trivial fixed-label coordinate]
\label{lem:perron-fixed-coordinate}
Assume \(\eta<\lambda\).  For every non-trivial irreducible \(\rho\), the
series
\begin{equation}\label{eq:perron-fixed-determinant}
  \mathfrak F_\rho:=
  \sum_{r,m\geq1}\frac{\mu(m)}{rm}
    \sum_{\substack{\pi\in\Irr(G)\\(r,m,\pi)\neq(1,1,\mathbf1)}}
       a_{\rho,\pi}^{(m)}
       \bigl(-\Log_{[0,1]}P_\pi(\lambda^{-rm})\bigr)
\end{equation}
is well defined and equals \(F_{\chi_\rho}(\lambda^{-1})\).  In
\eqref{eq:perron-fixed-determinant}, the notation means radial continuation
for the finitely many terms with \(rm=1\) and the normalized interior germ
for \(rm\geq2\).  The omitted scalar Perron triple has coefficient
\(a_{\rho,\mathbf1}^{(1)}=0\) and is deleted before endpoint evaluation.
\end{lemma}

\begin{proof}
Since \(\rho\neq\mathbf1\), character orthogonality gives
\(a_{\rho,\mathbf1}^{(1)}=0\).  If \(rm=1\), all remaining \(\pi\) are
non-trivial and \(P_\pi(t/\lambda)\neq0\) for \(0\leq t\leq1\) by the strict
gap.  For \(rm\geq2\), the trace expansion and
\(\abs{\Tr(A_\pi^n)}\leq(\dim A_\pi)\lambda^n\) give, uniformly for
\(0\leq t\leq1\),
\[
 \abs{-\Log_0P_\pi(t^{rm}\lambda^{-rm})}
 \leq C\sum_{n\geq1}\frac{\lambda^{(1-rm)n}}n
 \leq C'\lambda^{1-rm}.
\]
The Adams coefficients are uniformly bounded, and
\[
  \sum_{rm\geq2}\frac{\lambda^{1-rm}}{rm}
  =\sum_{k\geq2}\frac{d(k)}k\lambda^{1-k}<\infty.
\]
Dominated convergence in \eqref{eq:fixed-determinant-open-disc}, along
\(z=t/\lambda\), proves the assertion.
\end{proof}

\section{Boundary-point collisions and radial recovery}
\label{sec:boundary-collisions}

The full represented determinant family determines the primitive
length--class profile by \Cref{prop:comparison-dictionary}.  The vector of
class-product constants is substantially coarser: it retains only one radial
value of each non-trivial fixed-label coordinate.  We make this loss explicit
on a fixed base graph.

\begin{theorem}[Exact boundary collision on a fixed \(C_2\)-extension]
\label{thm:c2-boundary-collision}
Let \(G=C_2=\{e,s\}\), with non-trivial character
\(\varepsilon(s)=-1\), and let the base adjacency matrix be
\[
 A=\begin{pmatrix}
 0&4&0&0\\
 1&1&0&2\\
 0&0&0&4\\
 0&2&1&1
 \end{pmatrix}.
\]
There are two labellings \(\tau,\tau':E(\Gamma)\to C_2\) on the same marked
directed multigraph whose sign blocks are respectively
\[
 B=\begin{pmatrix}
 0&-2&0&0\\
 -1&-1&0&-2\\
 0&0&0&-2\\
 0&2&-1&1
 \end{pmatrix},
 \qquad
 B'=\begin{pmatrix}
 0&-4&0&0\\
 1&1&0&0\\
 0&0&0&-4\\
 0&0&1&1
 \end{pmatrix}.
\]
Both labellings satisfy the strict twisted gap, and all their class-product
constants agree:
\begin{equation}\label{eq:c2-collision-constants}
 K_{\{e\}}^\tau=K_{\{e\}}^{\tau'},
 \qquad
 K_{\{s\}}^\tau=K_{\{s\}}^{\tau'}.
\end{equation}
Nevertheless, their length-one primitive data are
\begin{equation}\label{eq:c2-collision-length-one}
 \bigl(p_{1,\{e\}}(\tau),p_{1,\{s\}}(\tau)\bigr)=(1,1),
 \qquad
 \bigl(p_{1,\{e\}}(\tau'),p_{1,\{s\}}(\tau')\bigr)=(2,0).
\end{equation}
In particular, equality of all Frobenius-class Mertens constants does not
determine the primitive length--class profile, even when the graph and group
are fixed and both extensions satisfy the strict gap.  The two labellings are
not related by vertex switching and global conjugacy.
\end{theorem}

\begin{proof}
Every row of \(A\) has sum four.  The underlying directed graph is strongly
connected and has loops, so \(A\) is primitive and \(\lambda=4\).  For each
entry, the numbers of \(e\)- and \(s\)-labelled edges are
\[
 \frac{A_{ij}+B_{ij}}2,\qquad \frac{A_{ij}-B_{ij}}2
\]
for \(\tau\), and the same formulas with \(B'\) for \(\tau'\).  Inspection
shows that these are non-negative integers.  Thus the displayed matrices are
realized on the same marked collection of parallel edges.

Put \(Q(z)=1-z+4z^2\).  Exact determinant calculation gives
\begin{align}
 \det(I-zB)&=1-z^2+4z^4=Q(z^2),
 \label{eq:c2-collision-first-determinant}\\
 \det(I-zB')&=1-2z+9z^2-8z^3+16z^4=Q(z)^2.
 \label{eq:c2-collision-second-determinant}
\end{align}
The characteristic polynomials are
\[
 \det(xI-B)=x^4-x^2+4,
 \qquad
 \det(xI-B')=(x^2-x+4)^2.
\]
For the first polynomial, the substitution \(y=x^2\) shows that
\(y\overline y=4\), hence every root has modulus \(\sqrt2\).  The roots of
\(x^2-x+4\) have modulus two.  Therefore
\begin{equation}\label{eq:c2-collision-gaps}
 \rad(B)=\sqrt2<4,
 \qquad
 \rad(B')=2<4.
\end{equation}

Write \(L_B(z)=-\log\det(I-zB)\) and define \(L_{B'}\) similarly, with the
germs normalized at zero.  In \(C_2\),
\(\psi^m\varepsilon=\varepsilon\) for odd \(m\) and
\(\psi^m\varepsilon=\mathbf1\) for even \(m\).  Since the trivial block is
the same for the two labellings, \eqref{eq:fixed-adams-mobius} gives, first on
\(|z|<1/4\),
\begin{align}
 F_{\varepsilon}^{\tau}(z)-F_{\varepsilon}^{\tau'}(z)
 &=\sum_{k\ge1}\frac1k
   \left(\sum_{\substack{m\mid k\\m\ \mathrm{odd}}}\mu(m)\right)
   \bigl(L_B(z^k)-L_{B'}(z^k)\bigr)\notag\\
 &=\sum_{j\ge0}2^{-j}
   \bigl(L_B(z^{2^j})-L_{B'}(z^{2^j})\bigr).
 \label{eq:c2-collision-dyadic-reduction}
\end{align}
Here the odd-divisor M\"obius sum is one when \(k\) is a power of two and
zero otherwise.  If \(q_j=\log Q(z^{2^j})\), then
\eqref{eq:c2-collision-first-determinant}--
\eqref{eq:c2-collision-second-determinant} turn the last series into
\[
 \sum_{j\ge0}2^{-j}(2q_j-q_{j+1})=2q_0.
\]
The series converges absolutely up to \(0\le z\le1/4\), since
\(q_j=O(z^{2^j})\).  Consequently,
\begin{equation}\label{eq:c2-collision-exact-difference}
 F_{\varepsilon}^{\tau}(z)-F_{\varepsilon}^{\tau'}(z)
 =2\log Q(z)\qquad(0\le z\le1/4).
\end{equation}
Because \(Q(1/4)=1\), the two sign coordinates agree at the Perron point.
The base matrix, and hence the scalar terms \(\gamma+\log C(A)\), are also
identical.  Formula \eqref{eq:correct-product-fourier} now proves
\eqref{eq:c2-collision-constants}.  For reference, direct reduction gives
\[
 \det(I-zA)=(1-4z)(1+2z-3z^2-4z^3),
 \qquad C(A)=\frac45.
\]

The only length-one primitive orbits are the loops at vertices two and four.
Their signs are respectively \((-1,1)\) for \(B\) and \((1,1)\) for \(B'\),
which proves \eqref{eq:c2-collision-length-one}.  Finally, vertex switching
acts on a sign block by diagonal sign conjugacy and leaves every loop label
unchanged.  Global conjugacy is trivial in \(C_2\).  Thus the two labellings
cannot be switching--conjugate.  Equivalently, their distinct characteristic
polynomials already preclude diagonal similarity.
\end{proof}

The telescoping in the proof is an exact kernel identity, rather than a
coincidence detected by truncation.

\begin{proposition}[Dyadic boundary functional]
\label{prop:dyadic-boundary-functional}
Let \(0<x<1\), and let \(Q\) be a polynomial with \(Q(0)=1\) and no zero on
the radial segment \([0,x]\).  Use the logarithm germ continued from zero
along that segment and define
\[
 \mathscr B_x(P):=-\sum_{j\ge0}2^{-j}\log P(x^{2^j})
\]
whenever the series converges.  Then
\begin{equation}\label{eq:dyadic-boundary-coboundary}
 \mathscr B_x\bigl(Q(z^2)\bigr)
 -\mathscr B_x\bigl(Q(z)^2\bigr)=2\log Q(x).
\end{equation}
Hence \(Q(x)=1\) places the non-constant determinant ratio
\(Q(z^2)/Q(z)^2\) in the kernel of this single boundary functional.
\end{proposition}

\begin{proof}
Put \(q_j=\log Q(x^{2^j})\).  Since \(q_j=O(x^{2^j})\), both series converge
absolutely, and shifting the first one gives
\[
 -\sum_{j\ge0}2^{-j}q_{j+1}
 +2\sum_{j\ge0}2^{-j}q_j=2q_0.
\]
This is \eqref{eq:dyadic-boundary-coboundary}.
\end{proof}

The polynomial identity above belongs to a slightly larger class of finite
rational certificates.  The positivity condition in the next statement is
essential: without it, the positive real powers defining the boundary product
need not exist.

\begin{proposition}[Positive rational Mahler certificates]
\label{prop:rational-mahler-certificate}
Let \(0<x<1\), let \(H\in\QQ(z)^\times\) satisfy \(H(0)=1\), and assume
that \(H(t)>0\) for \(0\le t\le x\).  Suppose that
\begin{equation}\label{eq:rational-mahler-coboundary}
 H(z)=\delta R(z):=\frac{R(z^2)}{R(z)^2}
\end{equation}
for some \(R\in\overline{\QQ}(z)^\times\) with \(R(0)=1\).  Then:
\begin{enumerate}[label=\textup{(\roman*)},leftmargin=*]
 \item the function \(R\) is unique and in fact belongs to
   \(\QQ(z)^\times\);
 \item it has neither a zero nor a pole on \([0,x]\), and \(R(t)>0\) on
   that interval;
 \item the positive real product
   \begin{equation}\label{eq:rational-mahler-product}
    \Pi_x(H):=\prod_{j\ge0}H(x^{2^j})^{2^{-j}}
   \end{equation}
   converges and satisfies
   \begin{equation}\label{eq:rational-mahler-telescope}
    \Pi_x(H)=R(x)^{-2}.
   \end{equation}
\end{enumerate}
Consequently, within the positive rational coboundary class
\[
 \mathscr R_x^+:=\{R\in\QQ(z)^\times:R(0)=1,
                         \ \delta R(t)>0\text{ on }[0,x]\},
\]
one has the exact kernel identity
\begin{equation}\label{eq:positive-rational-certificate-kernel}
 \ker(\Pi_x)\cap\delta\mathscr R_x^+
 =\{\delta R:R\in\mathscr R_x^+,\ R(x)=1\}.
\end{equation}
\end{proposition}

\begin{proof}
First, \(\delta\) is injective on rational functions normalized to have value
one at zero.  Indeed, if \(U(0)=1\) and \(U(z^2)=U(z)^2\), write
\(U(z)=1+az^n+O(z^{n+1})\) with \(a\ne0\) and \(n\) least.  The first
non-constant term has degree \(2n\) on the left and degree \(n\) on the
right, a contradiction.  Hence \(U=1\).

For every automorphism
\(\sigma\in\operatorname{Gal}(\overline{\QQ}/\QQ)\),
\eqref{eq:rational-mahler-coboundary} gives
\(\delta(\sigma R/R)=1\) and \((\sigma R/R)(0)=1\).  Injectivity gives
\(\sigma R=R\), so \(R\in\QQ(z)^\times\); it also proves uniqueness.

Let \(v_R(t_0)\) denote the zero or pole order of \(R\) at a real
\(t_0\in(0,x]\).  Since \(H(t_0)\) is finite and non-zero,
\eqref{eq:rational-mahler-coboundary} implies
\[
 v_R(t_0^2)=2v_R(t_0).
\]
If \(v_R(t_0)\ne0\), iteration produces infinitely many distinct zeros or
poles accumulating at zero, contrary to rationality and \(R(0)=1\).
Thus \(R\) has none on \([0,x]\); continuity and \(R(0)=1\) then give
\(R>0\) there.

Finally, \(\log H(t)=O(t)\) at zero, so
\eqref{eq:rational-mahler-product} converges absolutely after taking real
logarithms.  Its finite truncations telescope exactly:
\[
 \prod_{j=0}^{N}
 \left(\frac{R(x^{2^{j+1}})}{R(x^{2^j})^2}\right)^{2^{-j}}
 =\frac{R(x^{2^{N+1}})^{2^{-N}}}{R(x)^2}.
\]
The numerator tends to one because \(R(0)=1\), proving
\eqref{eq:rational-mahler-telescope} and then
\eqref{eq:positive-rational-certificate-kernel}.
\end{proof}

\begin{remark}[Boundary of the rational certificate]
\label{rem:mahler-kernel-open-interface}
\Cref{prop:rational-mahler-certificate} classifies only the intersection of
the boundary kernel with the positive rational coboundary image.  It does not
assert the converse lifting
\[
 H\in\QQ(z)^\times,\quad H>0\text{ on }[0,x],\quad \Pi_x(H)=1
 \quad\Longrightarrow\quad H=\delta R.
\]
The present paper leaves that implication open.  The lifting and purity
results of Adamczewski--Faverjon concern linear regular-singular Mahler
systems \cite{AdamczewskiFaverjon2018MahlerSeveralVariablesII}.  Greuel's
implicit-function theorem requires a strict degree inequality which is not
met by the boundary equation
\(Y(z)^2=H(z)^2Y(z^2)\) in its critical quadratic case
\cite{Greuel2000ImplicitMahler}.  Lam proves transcendence for a particular
weighted Mahler product at non-zero algebraic points, but not multiplicative
relations for arbitrary rational ratios \cite{Lam2025ValuationProducts}.
Thus none of these results supplies the omitted converse, and no such
converse is used below.

The realizability restriction does not remove the special-value question.
Indeed, fix \(r\ge2\), an odd integer \(a>0\), and a positive even integer
\(b\), and let \(A\) have diagonal entries \(a\) and off-diagonal entries
\(b\).  Then \(A\) is positive and has Perron root
\(\lambda=a+(r-1)b\).  For any odd integers \(c_i,d_i\) with
\(\lvert c_i\rvert,\lvert d_i\rvert\le a\), the diagonal matrices
\[
 B=\operatorname{diag}(c_1,\ldots,c_r),\qquad
 B'=\operatorname{diag}(d_1,\ldots,d_r)
\]
are compatible sign blocks on this same base and satisfy
\(\rad(B),\rad(B')\le a<\lambda\).  Their determinant ratio is
\[
 \frac{\det(I-zB)}{\det(I-zB')}
 =\frac{\prod_i(1-c_i z)}{\prod_i(1-d_i z)}.
\]
Thus the omitted converse already contains multiplicative-relation questions
for weighted products built from these linear factors; it is not an artifact
of dropping the same-base compatibility or strict-gap hypotheses.
\end{remark}

This mechanism is strictly weaker than full \(G\)-cospectrality, equivalently
represented cospectrality in every irreducible representation:
\eqref{eq:c2-collision-first-determinant} and
\eqref{eq:c2-collision-second-determinant} are different.  Represented
cospectrality in every irreducible representation instead requires equality
of the full twisted spectra and therefore preserves the complete periodic profile
\cite{BoyleSchmieding2017FiniteGroupExtensions,
CavaleriDonno2022Cospectrality}.

We next delimit the finite minimality assertion.  A \emph{regular polynomial
\(C_2\)-certificate} is a tuple \((A,B,B',Q)\) in which \(A\) is a primitive
\(m\)-out-regular integral adjacency matrix, \(B,B'\) are compatible sign
blocks, \(Q\in\ZZ[z]\) is non-constant with \(Q(0)=Q(1/m)=1\), both sign
blocks have spectral radius less than \(m\), and
\begin{equation}\label{eq:regular-polynomial-certificate}
 \det(I-zB)=Q(z^2),
 \qquad
 \det(I-zB')=Q(z)^2.
\end{equation}
Its complexity is ordered lexicographically as
\[
 \bigl(2,\deg Q,|V|,m,|E|\bigr).
\]

\begin{proposition}[Minimality in the regular polynomial certificate class]
\label{prop:c2-certificate-minimality}
The certificate in \Cref{thm:c2-boundary-collision} has complexity
\((2,2,4,4,16)\), and this is the least complexity of a regular polynomial
\(C_2\)-certificate.
\end{proposition}

\begin{proof}
The first coordinate is fixed by the non-trivial binary group.  Since
\(Q(0)=Q(1/m)=1\), a non-constant \(Q\) cannot have degree one, so
\(\deg Q\ge2\).  If \(\deg Q=2\), then \(Q(z^2)\) has degree four, whereas
\(\deg\det(I-zB)\le |V|\); hence \(|V|\ge4\).

It remains to exclude \(m=2,3\) at \(\deg Q=2\) and \(|V|=4\).  Write
\begin{equation}\label{eq:minimality-q-form}
 Q(z)=1-kz+kmz^2,
 \qquad k\in\ZZ\setminus\{0\},
\end{equation}
which follows from \(Q(1/m)=1\).  When \(m=2\), the characteristic
polynomial forced by the second identity in
\eqref{eq:regular-polynomial-certificate} is
\((x^2-kx+2k)^2\).  If its roots have modulus strictly less than two, their
product gives \(|k|<2\).  The choice \(k=-1\) has a root \(-2\), and
\(k=0\) makes \(Q\) constant; therefore strict gap forces \(k=1\).

There are exactly ten non-negative integral four-entry rows of sum two,
namely \(2e_i\) and \(e_i+e_j\) for \(i<j\).  We exhaust the resulting
\(10^4\) matrices.  Primitivity is decided exactly by testing positivity of
a power through exponent ten, the four-dimensional Wielandt bound.  This
leaves \(2208\) primitive base matrices.  For each such matrix, every
compatible signed entry is enumerated from
\[
 B_{ij}\in\{-A_{ij},-A_{ij}+2,\ldots,A_{ij}\}.
\]
For an integral \(4\times4\) matrix, put \(t_r=\Tr(B^r)\).  Newton's
identities give
\[
 \det(I-zB)=1-t_1z+\frac{t_1^2-t_2}{2}z^2
 -\frac{t_1^3-3t_1t_2+2t_3}{6}z^3+\det(B)z^4.
\]
Thus the two target polynomials
\(1-z^2+2z^4\) and \((1-z+2z^2)^2\) are tested by the integer quadruples
\[
 (t_1,t_2,t_3,\det B)=(0,2,0,2),
 \qquad (2,-6,-10,4),
\]
respectively.  The exhaustive integer enumeration finds \(48\) primitive
bases supporting the first target and no primitive base supporting the
second.  Hence no certificate exists for \(m=2\).  The enumeration and all
counts are reproduced by the exact-arithmetic verifier cited in
\Cref{sec:exactness}.

When \(m=3\), compatibility gives
\(B\mathbf1\equiv A\mathbf1\equiv\mathbf1\pmod2\).  Therefore
\(I-B\) is singular over \(\mathbb F_2\) and \(\det(I-B)\) is even.  On the
other hand, the first identity in
\eqref{eq:regular-polynomial-certificate} and
\eqref{eq:minimality-q-form} give
\[
 \det(I-B)=Q(1)=1-k+3k=1+2k,
\]
which is odd.  Thus \(m=3\) is impossible.  Since \(m>1\), it follows that
\(m\ge4\), and regularity gives \(|E|=m|V|\ge16\).  The matrices in
\Cref{thm:c2-boundary-collision} attain every bound.
\end{proof}

The minimality assertion is deliberately confined to the finite polynomial
certificate class in \eqref{eq:regular-polynomial-certificate}; it does not
purport to classify possible isolated equalities of transcendental boundary
values lacking such a certificate.

Although one Perron-point value loses information, an analytic radial profile
restores it.  The following is a sufficient recovery criterion; no optimality
claim about the amount of profile data is intended.

\begin{theorem}[Radial-profile recovery]
\label{thm:radial-profile-recovery}
Let \(\tau,\tau':E(\Gamma)\to G\) be two labellings on the same primitive
base matrix, both satisfying the strict twisted gap.  For every conjugacy
class \(C\), define
\begin{equation}\label{eq:radial-profile-difference}
 D_C(z):=\sum_{n\ge1}
 \bigl(p_{n,C}(\tau)-p_{n,C}(\tau')\bigr)\log(1-z^n),
\end{equation}
where the logarithm is normalized at zero.  There is a number
\(1<\beta<\lambda\) for which every \(D_C\) is analytic on
\(|z|<\beta^{-1}\).  If, for every \(C\), \(D_C\) vanishes on a set
\(S\subset\{|z|<\beta^{-1}\}\) having an accumulation point in that disc,
then
\begin{equation}\label{eq:radial-profile-recovery}
 p_{n,C}(\tau)=p_{n,C}(\tau')
 \qquad(n\ge1,\ C\in\operatorname{Conj}(G)).
\end{equation}
It is enough, for example, to take
\(S=\{\lambda^{-r}:r\ge2\}\).
\end{theorem}

\begin{proof}
Let \(\eta_\tau,\eta_{\tau'}<\lambda\) be the two strict-gap radii.  The
scalar primitive contribution cancels because the base matrix is fixed.
The character expansion \eqref{eq:class-indicator} and the bound
\eqref{eq:explicit-primitive-trace-bound}, applied to both labellings, show
that for any
\[
 \max\{\eta_\tau,\eta_{\tau'},\sqrt\lambda\}<\beta<\lambda
\]
one has
\[
 p_{n,C}(\tau)-p_{n,C}(\tau')=O(\beta^n/n).
\]
Consequently, \eqref{eq:radial-profile-difference} converges normally on
compact subsets of \(|z|<\beta^{-1}\), and hence defines an analytic
function there.

If \(D_C\) vanishes on \(S\), the analytic identity theorem gives
\(D_C\equiv0\).  Suppose that \(n_0\) is the least length at which the two
primitive counts for this class differ.  In the expansion
\[
 D_C(z)=-\sum_{n\ge1}\sum_{r\ge1}
  \frac{p_{n,C}(\tau)-p_{n,C}(\tau')}{r}z^{nr},
\]
every contribution to the coefficient of \(z^{n_0}\) from a proper divisor
of \(n_0\) vanishes by minimality.  That coefficient is therefore
\(-\bigl(p_{n_0,C}(\tau)-p_{n_0,C}(\tau')\bigr)\ne0\), a contradiction.
This proves \eqref{eq:radial-profile-recovery}.  Finally,
\(\lambda^{-r}<\beta^{-1}\) for \(r\ge2\), and these points accumulate at
zero inside the disc.
\end{proof}

For the collision in \Cref{thm:c2-boundary-collision},
\eqref{eq:c2-collision-exact-difference} equals \(2\log Q(z)\), and
\(Q(z)=1\) on \(0<z\le1/4\) only at \(z=1/4\).  Thus its equality is an
isolated Perron-boundary evaluation and cannot persist on any set allowed by
\Cref{thm:radial-profile-recovery}.

\section{Exact verification boundary}
\label{sec:exactness}

The proofs above are symbolic.  The accompanying exact-arithmetic program
\path{artifacts/verify_a5_results.py} independently checks the displayed
four-vertex determinants and spectral gaps, reconstructs
\(R=1-z+4z^2\), isolates \(1/4\) as its sole collision in
\((0,1/4]\), and verifies the degree, height, coefficient-recursion, and
Pad\'e identities used in the effective criterion.  It also exhausts the
four-vertex two-out-regular bases used in
\Cref{prop:c2-certificate-minimality}: there are \(2208\) primitive bases,
\(48\) support the first target determinant, and none supports the second.

The program records two proof-boundary audits.  It checks the parameter
specialization \(p=2,N=0,n=1,m=M=2,U=L=1\) and \(4<8\) in Nishioka's
theorem, but does not claim to machine-verify that published transcendence
theorem.  It also computes the \(C_3\) Adams--M\"obius state vectors through a
finite range and confirms non-zero coefficients at primes
\(2,5,11,17\); the infinitude assertion in the manuscript is the elementary
number-theoretic argument in the proof, not an inference from this finite
check.

The focused unit test is \path{artifacts/test_verify_a5_results.py}.  Exact
commands and expected final lines are listed in \path{REPRODUCE.md}.  These
computations supplement the proofs and are not used as substitutes for
analytic or parameter-uniform arguments.

\section{Conclusion}
\label{sec:conclusion}

For arbitrary binary extensions on one \(v\)-vertex base, fewer than
\(16v^3\) algebraic radial samples in the open Perron interval determine the
complete primitive length--class profile; the strict twisted gap is needed
only if the Perron endpoint is sampled.  The determinant ratio governing the
test is canonically the ratio of the standard Artin--Mazur zeta functions of
the two binary covers.  The proof is genuinely effective:
one collision produces a unique rational Mahler certificate, its numerator
and denominator satisfy explicit input-dependent degree and height bounds,
and a finite Pad\'e system reconstructs or rejects it.  The exact four-vertex
pair shows why a single Perron-boundary value cannot be an inverse invariant.

The strongest implication in this chain is the determinant lifting theorem.
Its interior form is unconditional, while its endpoint form uses the strict
gap only to continue through the Perron circle.  The arithmetic proof uses
more than positivity of a rational function:
same-base determinant parity supplies integral coefficients for the critical
Mahler product, Nishioka's theorem converts an algebraic special value into
algebraicity of the function, and the explicit algebraic-to-rational descent
then supplies the coboundary.  Removing the parity hypothesis would require a
different arithmetic input.

The same mechanism does not presently yield a polynomial sample bound for
every finite group.  For \(C_3\), Adams--M\"obius inversion already couples
two character determinants at infinitely many arithmetically distributed
powers, rather than one lacunary Mahler orbit.  A general finite-group theorem
therefore requires a new special-value and effective zero theory for these
coupled Euler systems.  In its absence, the proved general statement remains
unconditional analytic recovery from a profile set with an interior
accumulation point.


\appendix
\section{The determinant-derived radial boundary}
\label{app:radial-boundary}

The strict-gap theorem uses only zero-free radial continuation.  For
completeness, this appendix records the general determinant-derived convention
at a peripheral endpoint.

\begin{proposition}[Canonical radial boundary construction]
\label{prop:derived-boundary}
Let \(P_\pi(z)=\det(I-zA_\pi)\).  The normalized germ
\(-\Log_0P_\pi(z)\), its continuation along \(z=t/\lambda\) with
\(0\leq t<1\), its peripheral factorization, and every radial renormalized
endpoint are determined by \(P_\pi\), \(\lambda\), and this path convention.
More explicitly, factor
\begin{equation}\label{eq:peripheral-factorization}
  P_\pi(t/\lambda)
   =\prod_{\zeta\in\Omega_\pi}(1-t\zeta)^{m_{\pi,\zeta}}
      Q_{\partial,\pi}(t),
\end{equation}
where \(\Omega_\pi\subset\TT\) is the multiset of peripheral phases and
\(Q_{\partial,\pi}\) has no zero on \([0,1]\).  If
\(s_\pi=m_{\pi,1}\), then
\begin{align}
  B_{1,\pi}^{\mathrm{ren}}
  &:={\lim}_{t\uparrow1}
     \left(-\Log_{[0,t]}P_\pi(t/\lambda)+s_\pi\log(1-t)\right)
     \label{eq:derived-b-ren}\\
  &=-\sum_{\substack{\zeta\in\Omega_\pi\\\zeta\neq1}}
       m_{\pi,\zeta}\Log_{[0,1]}(1-\zeta)
     -\Log_{[0,1]}Q_{\partial,\pi}(1).
     \label{eq:derived-b-factor}
\end{align}
Thus the branch, peripheral factors, and renormalized value are derived from
the determinant polynomial rather than supplied as independent cocycle data.
\end{proposition}

\begin{proof}
The roots and their multiplicities in \eqref{eq:peripheral-factorization}
are determined by \(P_\pi(t/\lambda)\).  The normalized logarithm has a unique
continuation along each zero-free subinterval beginning at zero.  Removing
\((1-t)^{s_\pi}\), taking the radial limit, and retaining that continuation
on every other factor gives \eqref{eq:derived-b-factor}.
\end{proof}

\section{Periodic data, gain-graph cospectrality, and fixed presentations}
\label{sec:inverse-background}

This section is included to delimit, rather than enlarge, the claims of the
paper.  Group-algebra adjacency matrices, represented spectra, switching,
and \(G\)-cospectrality are developed systematically in
\cite{CavaleriDAngeliDonno2021BalanceGainGraphs,
CavaleriDonno2022Cospectrality}; the foundational gain-graph and voltage-graph
settings are \cite{Zaslavsky1989BiasedGraphs,
GrossTucker1987TopologicalGraphTheory}.  Boyle and Schmieding's
conjugacy-class periodic data are a complete, finitely computable invariant of
the extension restricted to periodic points
\cite{BoyleSchmieding2017FiniteGroupExtensions}.  The theorem below merely
puts the determinant notation of the present paper into that established
dictionary.

\subsection{The periodic-data dictionary}

Let \(\CC[G]\) have basis \([g]\), \(g\in G\), and define the group-algebra
adjacency matrix
\begin{equation}\label{eq:group-algebra-adjacency}
  \mathbb A_\tau:=
  \left(\sum_{e:i\to j}[\tau(e)]\right)_{i,j\in V}
  \in M_V(\CC[G]).
\end{equation}
For \(x=\sum_ga_g[g]\in\CC[G]\), let
\[
  \kappa(x):=
  \left(\sum_{g\in C}a_g\right)_{C\in\operatorname{Conj}(G)}.
\]
Thus
\begin{equation}\label{eq:kappa-periodic-data}
  \kappa\!\left(\operatorname{tr}(\mathbb A_\tau^n)\right)
  =(N_{n,C}(\tau))_{C\in\operatorname{Conj}(G)}.
\end{equation}
We call equality of the vectors in \eqref{eq:kappa-periodic-data} for every
\(n\) \emph{directed \(G\)-cospectrality}.  This is the directed-multigraph
counterpart of the gain-graph notion; for finite groups, Cavaleri and Donno
prove in their setting that \(G\)-cospectrality is equivalent to cospectrality
in every unitary representation
\cite{CavaleriDonno2022Cospectrality}.

\begin{theorem}[Comparison of determinant, cospectral, and periodic data]
\label{thm:comparison-dictionary}
For two edge labellings \(\tau,\tau':E(\Gamma)\to G\) on the same fixed
marked directed multigraph, the following are equivalent:
\begin{enumerate}[label=\textup{(\roman*)},leftmargin=*]
  \item \(P_{\rho,\tau}(z)=P_{\rho,\tau'}(z)\) for every
    \(\rho\in\Irr(G)\);
  \item
    \(\Tr(A_{\rho,\tau}^n)=\Tr(A_{\rho,\tau'}^n)\) for every
    \(\rho\in\Irr(G)\) and \(n\geq1\);
  \item \(\tau\) and \(\tau'\) are directed \(G\)-cospectral:
    \[
      \kappa(\operatorname{tr}(\mathbb A_\tau^n))
      =\kappa(\operatorname{tr}(\mathbb A_{\tau'}^n))
      \quad(n\geq1);
    \]
  \item \(N_{n,C}(\tau)=N_{n,C}(\tau')\) for every \(n\geq1\) and every
    conjugacy class \(C\);
  \item \(p_{n,C}(\tau)=p_{n,C}(\tau')\) for every \(n\geq1\) and every
    conjugacy class \(C\).
\end{enumerate}
In particular,
\[
 \{P_{\rho,\tau}\}_{\rho\in\Irr(G)}
 \Longleftrightarrow
 \left\{\kappa(\operatorname{tr}(\mathbb A_\tau^n))\right\}_{n\geq1}
 \Longleftrightarrow
 \{N_{n,C}(\tau)\}_{n,C}.
\]
This is an equivalence of unmarked periodic data, not a new conjugacy or
switching classification.
\end{theorem}

\begin{proof}
The normalized logarithmic germ of \(P_{\rho,\tau}\) is
\[
 -\Log_0P_{\rho,\tau}(z)
 =\sum_{n\geq1}\Tr(A_{\rho,\tau}^n)z^n/n.
\]
Hence (i) and (ii) are equivalent.  Expanding
\(\operatorname{tr}(\mathbb A_\tau^n)\) over closed paths and then applying
\(\rho\) gives
\begin{equation}\label{eq:periodic-character-transform}
  \Tr(A_{\rho,\tau}^n)
   =\sum_C N_{n,C}(\tau)\chi_\rho(C).
\end{equation}
The irreducible character table is invertible on class functions, so
\eqref{eq:periodic-character-transform} proves the equivalence of
(ii)--(iv), consistently with \eqref{eq:kappa-periodic-data}.

Finally, decomposition of a period-\(n\) point into a primitive orbit of
length \(d\mid n\) and its \(n/d\)-fold traversal gives
\begin{equation}\label{eq:periodic-primitive-triangular}
  N_{n,C}
  =\sum_{d\mid n}d
     \sum_{\substack{D\in\operatorname{Conj}(G)\\D^{\,n/d}=C}}
       p_{d,D}.
\end{equation}
The term \(d=n\) is \(np_{n,C}\); every other term uses primitive lengths
strictly smaller than \(n\).  Thus \eqref{eq:periodic-primitive-triangular}
is triangular and recursively invertible, proving (iv)\(\Leftrightarrow\)(v).
\end{proof}

\begin{remark}[Relation to complete periodic data]
The data in (iv) are exactly the conjugacy-class periodic data emphasized by
Boyle--Schmieding, up to the choice between point and orbit normalization.
Their work establishes the relevant completeness on periodic points and
finite computability.  The passage to (i) is finite character Fourier
transform, while the passage to (v) is the triangular
Adams--M\"obius bookkeeping in \eqref{eq:periodic-primitive-triangular}.
Neither passage is presented here as an independent rigidity theorem.
\end{remark}

\subsection{Switching and a finite fundamental-cycle certificate}

A vertex switching function \(u:V\to G\) acts by
\begin{equation}\label{eq:switching-action}
  \tau^u(e)=u(o(e))^{-1}\tau(e)u(t(e)).
\end{equation}
It conjugates every twisted matrix by
\(\operatorname{diag}_{v\in V}\rho(u(v))\), and therefore implies every
condition in \Cref{thm:comparison-dictionary}.  The converse is false in
general; gain-graph cospectrality is deliberately weaker than switching
equivalence.  See \cite{CavaleriDonno2022Cospectrality} for the representation
theory and switching-class count, and
\cite{AbiadBelardoKhramova2024Switching} for constructions of cospectral gain
graphs by switching methods.

\begin{proposition}[Finite fundamental-cycle switching certificate]
\label{prop:fundamental-cycle-certificate}
Assume the underlying undirected graph of \(\Gamma\) is connected.  Choose a
spanning tree \(T\), a root, and one actual marked edge above every tree edge.
Switch both cocycles so that these chosen tree-edge labels are the identity.
Let \(e_1,\ldots,e_r\) be the remaining marked edges, where
\(r=|E|-|V|+1\), and let \(g_i,g_i'\) be the gains of the corresponding
oriented fundamental cycles for the two normalized cocycles.  Then
\(\tau\) and \(\tau'\) are switching equivalent if and only if there is one
\(h\in G\) such that
\begin{equation}\label{eq:simultaneous-fundamental-conjugacy}
  g_i'=h^{-1}g_i h\qquad(1\leq i\leq r).
\end{equation}
Thus switching is certified by finitely many fundamental cycles; no
spectral-projector compatibility and no quantification over all closed paths
is needed.
\end{proposition}

\begin{proof}
Tree normalization is obtained recursively from the root using
\eqref{eq:switching-action}; it is the usual voltage-graph normalization
\cite{GrossTucker1987TopologicalGraphTheory}.  Once every tree edge is
normalized, a cocycle is determined by the chord gains
\((g_1,\ldots,g_r)\), equivalently by a homomorphism from the free group
\(\pi_1(|\Gamma|)\) to \(G\).  The only switching freedom preserving the tree
normalization is the common value \(h\) at every vertex, and it acts by
simultaneous conjugation.  This proves \eqref{eq:simultaneous-fundamental-conjugacy};
it is also the standard gain-graph switching criterion
\cite{Zaslavsky1989BiasedGraphs,CavaleriDAngeliDonno2021BalanceGainGraphs,
CavaleriDonno2022Cospectrality}.
\end{proof}

\subsection{Presentation-relative determinant fibres}

Fix the graph, every vertex, and every edge name.  Write
\[
  H^1(\Gamma;G):=G^{E(\Gamma)}/G^{V(\Gamma)}
\]
for the switching quotient under \eqref{eq:switching-action}, and define
\[
  S_{\Gamma,G}([\tau])
  :=\bigl(P_{\rho,\tau}(z)\bigr)_{\rho\in\Irr(G)}.
\]

\begin{definition}[Presentation-relative determinant-fibre multiplicity]
\label{def:presentation-relative-mu}
For a fixed marked-edge presentation \(\Gamma\), set
\begin{equation}\label{eq:presentation-relative-mu}
  \mu_{\Gamma,G}([\tau])
  :=\#\left\{[\alpha]\in H^1(\Gamma;G):
     S_{\Gamma,G}([\alpha])=S_{\Gamma,G}([\tau])\right\}.
\end{equation}
\end{definition}

The argument \([\tau]\) will not be suppressed.  By definition,
\(\mu_{\Gamma,G}([\tau])=1\) if and only if this particular determinant
fibre contains one switching class.  This is a fibre-cardinality criterion,
not an intrinsic obstruction theorem.  A marked graph isomorphism transports
\eqref{eq:presentation-relative-mu} bijectively and hence preserves its
value.  No invariance under state splitting, higher-block recoding, or an
arbitrary automorphism of the base shift is asserted.

\begin{proposition}[What the bouquet multiplicity counts]
\label{prop:bouquet-presentation-relative}
Let \(\mathcal B_m\) be the one-vertex graph with \(m\) named loops and let
\(G\) be finite abelian.  Put
\[
  n_g(\tau):=\#\{e\in E(\mathcal B_m):\tau(e)=g\}.
\]
Then
\begin{equation}\label{eq:bouquet-named-fibre}
  \mu_{\mathcal B_m,G}([\tau])
   =\frac{m!}{\prod_{g\in G}n_g(\tau)!}.
\end{equation}
This number counts permutations of named loops.  The symmetric group \(S_m\)
acts transitively on that fibre, so its quotient by loop permutations has
cardinality one.
\end{proposition}

\begin{proof}
For a character \(\chi\in\widehat G\), the twisted matrix is the scalar
\[
  A_\chi(\tau)=\sum_{e\in E(\mathcal B_m)}\chi(\tau(e))
  =\sum_{g\in G}n_g(\tau)\chi(g).
\]
Finite Fourier inversion recovers the multiplicity vector \((n_g)_{g\in G}\)
from these scalars.  Thus the determinant fibre consists precisely of the
assignments of this multiset to the \(m\) named loops, which proves
\eqref{eq:bouquet-named-fibre}.  Loop permutations act transitively on those
assignments.
\end{proof}

\begin{corollary}[Qualification of the two-loop example]
\label{cor:two-loop-qualification}
For \(G=\ZZ/2\ZZ\), the two labelled bouquets
\[
  (\tau(a),\tau(b))=(0,1),
  \qquad
  (\tau'(a),\tau'(b))=(1,0)
\]
are distinct in the fixed named-edge switching quotient and have the same
determinant family.  Nevertheless, the symbol transposition \(s(a)=b\),
\(s(b)=a\) satisfies \(\tau'=\tau\circ s^{-1}\), and
\[
  (x,g)\longmapsto(sx,g)
\]
conjugates the two skew products.  Hence this is only the smallest
\emph{fixed-named-edge} determinant collision; it is not a smallest pair of
nonconjugate extensions.
\end{corollary}

\begin{proof}
The determinant assertion is the case \(m=2\) of
\Cref{prop:bouquet-presentation-relative}.  The displayed map commutes with
the shifts and intertwines the cocycles because it exchanges the two edge
symbols.
\end{proof}

\begin{remark}[Automorphism-quotiented inverse questions]
A genuinely dynamical inverse fibre would have to be formed after an action
such as
\[
  \Aut(X_\Gamma^+,\sigma)\backslash H^1(\Gamma;G),
\]
with care about which automorphisms preserve the chosen cocycle category.
No automorphism-free collision is supplied here, and no novelty is claimed
for inverse impossibility beyond the established gain-graph and periodic-data
literature.  This is why the inverse material is background rather than a
headline contribution.
\end{remark}


\section*{Competing interests}
The authors declare none.

\section*{Data and code availability}
No research datasets were generated or analysed. The separately submitted
Supplementary Material and reproducibility archive contain the companion
results, exact verification scripts, unit tests, and readable archived
outputs.

\section*{Guide to the references}
The trace and dynamical-zeta formalism is taken from
\cite{BowenLanford1970Zeta,ParryPollicott1990Zeta}; necklace/Witt inversion
is taken from \cite{MetropolisRota1983Necklaces}.  The arithmetic inputs are
Kumiko Nishioka's verified special-value theorem, the algebraic-solution
rationality statement attributed to Keiji Nishioka and quoted as the
unverified assumption \textup{(KN85)} where needed, and the
Dieudonn\'e--Dwork criterion
\cite{Nishioka1982MahlerFunctionValues,Nishioka1985AlgebraicSolutions,
PomeratStraub2024RootsPowerSeries}.  Direct Frobenius-class antecedents are
\cite{ParryPollicott1986Chebotarev,NooraniParry1992ChebotarevShifts,
MeslizaNoorani1999FrobeniusMertens,Noorani2003ClosedOrbitsToral}.
The nearest algorithmic comparisons are
\cite{ChyzakDreyfusDumasMezzarobba2018MahlerSolutions,
ChyzakDreyfusDumasMezzarobba2025FirstOrderFactors}.  The inverse-data
appendix uses Boyle--Schmieding's conjugacy-class periodic data and the
gain-graph literature
\cite{BoyleSchmieding2017FiniteGroupExtensions,
CavaleriDAngeliDonno2021BalanceGainGraphs,CavaleriDonno2022Cospectrality,
AbiadBelardoKhramova2024Switching,Zaslavsky1989BiasedGraphs,
GrossTucker1987TopologicalGraphTheory}.

\bibliographystyle{amsplain}
\bibliography{references}

\end{document}
